% !TEX root = ../paper.tex

\section{Data and Sample Construction}
\label{app:data_sources}
\label{sec:appendix_data}

This appendix documents the data sources and sample construction procedures used in the analysis.

\subsection{Data Sources}

The analysis draws on the following administrative and proprietary datasets:

\paragraph{Cook County Assessor Data.}
Sales records, property assessments, and improvement characteristics for residential parcels in Cook County.

\paragraph{RentHub Listings.}
RentHub rental listings for Illinois from 2014--2022, restricted to Chicago using the ward assignment procedure described below.

\paragraph{Chicago Building Permits.}
Permit records and processing times for project types where aldermen are most likely to matter.

\paragraph{Ward Boundary Shapefiles.}
Official ward boundaries for the two eras used in the paper: pre-2015 and 2015--2023.

\paragraph{Amenity Data.}
School locations, Chicago Park District park-boundary polygons, City of Chicago major streets, CTA stops, and Lake Michigan geometry used to construct distance controls for the rent and sales RDs.
CTA distances use the permanent station network in place during the listing month or on the sale date.

\paragraph{Census Data.}
American Community Survey block group demographics and 2010/2020 decennial census block boundaries.
I construct ward-level demographic controls by overlaying block groups with the ward boundaries in effect in each year.
Population, household, tenure, race, ethnicity, and education counts are allocated to wards in proportion to each block group's area in the ward, and ward shares are calculated from those allocated counts.
Ward median household income is the household-weighted mean of block group medians using the same allocation.
Where historical ward polygons overlap slightly, I scale the ward-specific intersections to the block group's area inside the union of Chicago wards so that no area is counted twice.

\subsection{High-Discretion Permits: Sample Construction}
\label{subsec:permits_sample}

The building permits data come from the City of Chicago's open data portal and contain permits recorded by the Department of Buildings since 2006.

\paragraph{Sample Restrictions.}
I construct the high-discretion sample using the permit types defined in Section~\ref{subsec:permit_pipeline}.

All remaining permit types are excluded from the high-discretion sample.
In placebo analyses, low-discretion permits are the remaining routine permit types, excluding signs.
The event-study outcomes retain observed permit records except those marked cancelled, revoked, or suspended.

\paragraph{Processing Time.}
For each permit, I compute processing time as the number of days between the application date and the issue date.
Permits with missing dates or negative processing times are excluded.
Same-day permits enter the lagged permit-volume control but not the log processing-time regression.

\paragraph{Ward Assignment.}
Permits are assigned to wards based on the permit's geographic coordinates and the ward boundaries in effect at the time of application.

\subsection{Permit Event-Study Panel Construction}
\label{subsec:event_study_panel}

For the permit event study, I construct a block-level treatment panel around the 2015 redistricting.

\paragraph{Treatment Assignment.}
Blocks are classified as treated if they were redistricted from one ward to another.
In the 2015 cohort, I compare ward assignments under the pre-2015 map (2014 boundaries) to post-2015 boundaries.
Census blocks are assigned to wards by largest-area overlap after projecting blocks and wards to EPSG:3435.

\paragraph{Predetermined Stringency.}
Treatment assignment is based on the difference between the destination and origin aldermen's baseline stringency scores.
For this design only, I re-estimate the stringency scores using permits filed from 2006 through 2014, before the redistricting took effect.
Origin and destination aldermen are assigned using the 2014 pre-redistricting ward-to-alderman mapping.
The change in stringency for block $i$ is
\[
\Delta\text{Stringency}_i = \text{Stringency}_{\text{dest}} - \text{Stringency}_{\text{origin}}.
\]
Blocks with a positive difference are classified as assigned toward greater stringency, blocks with a negative difference as assigned toward greater leniency, and blocks that remain in the same ward form the comparison group.

\paragraph{Assigned Stringency.}
For each block, I measure the change in stringency using the 2014 incumbents in its origin and destination wards.
Some wards elected new aldermen when the new map took effect, so this assigned change does not always match the change under the alderman who ultimately took office.
I nevertheless keep treatment tied to the pre-election incumbents.
Defining treatment using the 2015 winner could mix the effect of aldermanic stringency with the ward's own demand for more or less development.
I retain all blocks with complete pre- and post-map ward assignments regardless of electoral turnover and interpret the estimate as the effect of assigned stringency.

\paragraph{Analysis Sample.}
Permits are aggregated to block-year panels using application year.
The event study uses the 2015 implementation cohort with relative years $-5$ through $5$.
The 500ft sample restriction is based on the distance from each census block polygon's centroid to the cohort's pre-redistricting boundary geometry.
Controls use the nearest incident pre-redistricting ward boundary, while treated blocks use the fixed origin-destination ward boundary created by redistricting.
The regressions control for each block's total permit volume in years $-5$ through $-1$ and an indicator for having no permits in those years, with both controls interacted with calendar year.

\subsection{New Construction: Sample Construction}
\label{subsec:new_construction_sample}

The new-construction sample combines residential improvement records and commercial valuation records from the Cook County Assessor.
I reconcile the two sources at the project level so that predecessor and successor PINs, split parcels, and duplicate representations of the same construction episode do not enter as separate projects.
I also compare the resulting ledger with every issued residential new-construction permit.
This recovers completed projects absent from the initial Assessor ledger and identifies permit records that cannot support comparable whole-building density measures.

\paragraph{Sample Restrictions.}
I apply the following restrictions: (1) lot area is positive; (2) building area is positive; (3) at least one dwelling unit; and (4) construction year falls between 2006 and 2022.
These restrictions exclude vacant parcels, data errors, and buildings outside the active sample window.
Projects with a completed permit but no defensible combination of whole-building units, building area, and lot area are not added because their density outcomes would not be comparable to the Assessor-based measures.

\paragraph{Construction Year.}
I use the selected Assessor construction episode and retain permit or completion evidence when it is needed to resolve a project whose reported year changes across records.
For multiple-card parcels, the project ledger suppresses successor records that represent the same construction and consolidates prorated PIN groups whose cards repeat whole-building characteristics.
The remaining cases that cannot be resolved by the source-priority rules are recorded in a small, frozen decision ledger.

\paragraph{Geocoding.}
I first locate each original PIN using the Historical Parcel Universe for the construction year.
When the historical service has no matching parcel, I use an exact current-PIN centroid or an accepted historical address match.
The project ledger retains the source of each location and suppresses predecessor or successor parcels that would duplicate another construction project.

\paragraph{Multifamily Classification.}
Assessor class is the primary programmatic signal for distinguishing multifamily projects from collections of single-family buildings.
Permit descriptions and external property records recover unit counts when the Assessor value is not a credible whole-building count.
A frozen exception ledger handles the small number of projects for which the general class and permit rules do not determine a unique classification.

\paragraph{Zoning.}
I assign each project the broad zoning group in force on June 15 of its construction year.
The zoning history combines archived City zoning maps with passed map amendments and uses historical parcel boundaries to follow PIN changes over time.

\paragraph{Density Measures.}
For each parcel, I compute two density measures: Floor Area Ratio (FAR), defined as total building area divided by lot area; and Dwelling Units Per Acre (DUPAC), defined as the number of dwelling units divided by lot acreage.

\paragraph{Ward and Boundary Assignment.}
Each parcel is assigned to a ward based on its centroid location and the ward boundaries in effect at the time of construction.
Because the Assessor reports only the year built, I assign June 15 of that year as the construction date and use the ward map and alderman in office on that date.
Distance to the nearest ward boundary is computed after projecting parcels and boundaries to the Illinois East State Plane coordinate system and is reported in feet.

\subsection{Summary Statistics: New Construction Sample}

Table~\ref{tab:summary_stats} presents summary statistics for the exact 500ft common sample used in the main FAR and DUPAC regressions.
The sample contains 3,692 projects, including 822 classified as multifamily.
Column (1) reports all new construction and Column (2) reports the multifamily sample.

\subsection{Home Sales: Sample Construction}
\label{subsec:sales_sample}

The home sales data come from the Cook County Assessor's Office and include all residential property transactions in Chicago from 1999--2022.

\paragraph{Sample Restrictions.}
I apply the following restrictions to construct the analysis sample: (1) residential property classes only, including single-family homes (classes 202--210) and small multifamily buildings with 2--6 units (class 211), excluding mixed-use properties (class 212) and condominiums (classes 299, 399); (2) sale price exceeds \$10,000 to exclude nominal transfers and gifts; (3) market transactions only, defined as warranty or trustee deeds; (4) excludes land-only sales; (5) seller and buyer names are recorded and non-identical; and (6) single-parcel transactions only.

\paragraph{Price Winsorization.}
Sale prices are first deflated to 2022 dollars using the monthly Chicago CPI-U all-items series and then winsorized at the 1st and 99th percentiles of the 2006--2022 analysis sample to limit the influence of outliers on mean calculations.

\paragraph{Geocoding.}
Each sale is geocoded by matching its full 14-digit property identification number (PIN) to the Cook County historical parcel universe in the sale year.
If an exact sale-year coordinate is unavailable, I use the PIN's coordinate from the 2025 parcel universe.

\paragraph{Ward Assignment.}
Sales are assigned to wards and to the aldermen in office based on the exact sale date.
For the paper's sample window, sales before May 18, 2015 use pre-2015 boundaries and sales on or after May 18, 2015 use the 2015 ward boundaries.

\paragraph{Distance to Ward Boundary.}
Distance calculations follow the same procedure as for new construction (Section~\ref{subsec:new_construction_sample}).
I project the data to the Illinois East State Plane coordinate system (EPSG:3435) before computing signed distance based on relative aldermanic stringency.

\subsection{RentHub Listings: Sample Construction}
\label{subsec:renthub_sample}

The rental panel starts from scrape-level listing rows.
I clean addresses and listing characteristics before constructing proxy identifiers.
Address strings coded as zero, blank, NA, N/A, NULL, NONE, or UNKNOWN are treated as missing.
Building types are recoded with townhouse and condo categories handled before broader house or commercial categories.

Because the raw data do not identify true units, I use a transparent floorplan proxy.
A property proxy combines normalized address and rounded coordinates; if address is missing, the proxy uses coordinates only and carries an address-missing flag.
A floorplan proxy adds beds, baths, square feet, and cleaned building type.
I first collapse raw rows to floorplan-by-scrape-date observations using median rent, then collapse to one floorplan-month observation using the median daily rent observed in that month.
This construction prevents daily scrape persistence from overweighting a listing while retaining multiple floorplans in large multifamily buildings.
When multiple rents appear for the same floorplan proxy on the same scrape date, I retain the group and use the median rent, while carrying flags for unusually large within-day rent dispersion.
Before aggregating to months, I trim rent and square footage at the 1st and 99th percentiles within year-by-bedroom cells.

I keep observations with positive rent and valid coordinates that can be assigned to Chicago wards.
Ward boundaries and aldermen are assigned using the listing's first observed date, with the May 2015 transition split at May 18.
The main rent analysis further restricts the sample to listings with reliable geographic assignments.
I exclude manually rejected locations, addresses with unstable coordinate clusters, and addressless listings concentrated at generic coordinates.
I also exclude listings when the modal coordinate for the address would change the assigned ward, neighboring ward, or ward pair, or would change the measured distance to the boundary by more than 100ft.
The regression requires nonmissing square footage, bathrooms, and amenity controls.
Bedrooms enter as categorical controls, so studios remain in the sample.
